% Part: first-order-logic
% Chapter: natural-deduction
% Section: soundness-identity

\documentclass[../../../include/open-logic-section]{subfiles}

\begin{document}

\olfileid{fol}{ntd}{sid}

\olsection{Kasahihan mawa \usetoken{S}{identity}}

\begin{prop}
Deduksi natural mawa aturan kanggo $\eq$ iku sah.
\end{prop}

\begin{proof}
Saben !!{formula} awujud $\eq[t][t]$ iku valid, amarga kanggo saben
!!{structure}~$\Struct M$, $\Sat{M}{\eq[t][t]}$. (Cathet yen kita
nganggep term $t$ iku tertutup, yaiku ora ngemot variabel, mula
penetapan variabel ora relevan).

Upamana inferensi pungkasan ing !!a{derivation} iku \Elim{\eq},
yaiku !!{derivation} kasebut awujud:
\begin{prooftree}
  \AxiomC{$\Gamma_1$}
  \RightLabel{$\delta_1$}
  \DeduceC{$\eq[t_1][t_2]$}
  \AxiomC{$\Gamma_2$}
  \RightLabel{$\delta_2$}
  \DeduceC{$!A(t_1)$}
  \RightLabel{\Elim{\eq}}
  \BinaryInfC{$!A(t_2)$}
\end{prooftree}
Premis $\eq[t_1][t_2]$ lan $!A(t_1)$ miturut urutane !!{derive}d
saka asumsi !!{undischarged}~$\Gamma_1$ lan $\Gamma_2$. Kita arep
nuduhake yen $!A(t_2)$ ndherek saka $\Gamma_1 \cup \Gamma_2$. Nimbang
!!a{structure}~$\Struct{M}$ kanthi $\Sat{M}{\Gamma_1 \cup
  \Gamma_2}$. Miturut hipotesis induksi, $\Sat{M}{!A(t_1)}$ lan
$\Sat{M}{\eq[t_1][t_2]}$. Mula, $\Value{t_1}{M} = \Value{t_2}{M}$.
Ayo $s$ dadi sembarang penetapan variabel, lan $m = \Value{t_1}{M} =
\Value{t_2}{M}$. Miturut
\olref[fol][syn][ext]{prop:ext-formulas}, $\Sat{M}{!A(t_1)}[s]$ yen
lan mung yen $\Sat{M}{!A(x)}[\Subst{s}{m}{x}]$ yen lan mung yen
$\Sat{M}{!A(t_2)}[s]$. Amarga $\Sat{M}{!A(t_1)}$, kita duwe
$\Sat{M}{!A(t_2)}$.
\end{proof}

\end{document}
